This thesis set out to work toward generative design of tuned mass resonator layouts on a plate: given a target equivalent radiated power spectrum, propose a resonator design that achieves it. To reach that goal, a design space of five variables per resonator, mass, tuning frequency, position and a damping ratio, was defined and sampled evenly with Latin hypercube sampling, and datasets of 100{,}000 two-resonator and 100{,}000 three-resonator configurations were generated with a modal physics solver. Ten forward neural network architectures were then built and compared as fast surrogates for that solver, and seven invertible architectures were developed and evaluated that predict a spectrum from a design and propose a design from a spectrum with a single set of weights.

The forward comparison showed that an explicit, permutation-invariant resonator encoder combined with physics-aware frequency features matters: every architecture that uses them clearly outperforms the plain feed-forward baseline, and the best architecture, DCO, reaches about 1.1\,dB root-mean-square error against the solver on unseen two-resonator designs. The top and the bottom of the ranking held on a three-resonator dataset, where the error of the architectures increased by a factor of 1.1 to 2.0 (DCO from 1.1\,dB to 2.2\,dB), confirming that the comparison reflects a real property of each architecture rather than an artefact of the simpler two-resonator case. A separate experiment, reported in Appendix~A, showed that these forward models interpolate well within the frequency band they are trained on but lose accuracy once asked to extrapolate beyond it, which is a useful caution for how the surrogates should be used.

In the forward direction the invertible architectures reached 2.4 to 3.8\,dB RMSE (2.7\,dB for iDON), which is comparable to the plain forward architectures on the harder three-resonator dataset but behind them on the two-resonator dataset (1.1 to 2.1\,dB for eight of the ten); since the datasets differ, this comparison is only indicative, but it suggests that making a core invertible costs some forward accuracy. Their inverse direction already recovers the overall shape of a target spectrum and its dominant tuning frequencies reasonably well, with correlations around 0.9 for the tuning frequency. Recovering exact resonator positions is markedly harder, with correlations as low as 0.12 for some architectures, because many different layouts can produce a nearly identical spectrum, as the non-uniqueness example in Chapter~\ref{ch:theory} showed directly. For iDON, padding the invertible branch beyond the design's own size made a clear difference on every metric checked, including halving the median position-recovery error, which shows that the capacity of the invertible core, not just its architecture, matters for how well the inverse direction performs.

These invertible results come from an earlier, four-variable version of the design space, built before the damping ratio and before the final two-resonator, 100{,}000-configuration forward dataset used in Chapter~\ref{ch:results}. Retraining the invertible architectures on the current five-variable design space is therefore the most direct next step: it is what would turn the forward surrogates and invertible architectures built in this thesis into one consistent generative design tool. Beyond that, four further steps follow from the results in this thesis. First, improving position recovery specifically, for example by giving the inverse direction an explicit way to represent several equally valid proposals rather than averaging over them, as the probabilistic inverse models of \textcite{thorpe2025generalizable} or the generative latent-space formulation of \textcite{zang2026designgenno} do. Second, extending the invertible models from two resonators to three, the way the forward comparison was already extended in Chapter~\ref{ch:results}. Third, applying the Monte Carlo uncertainty quantification procedure described in Chapter~\ref{ch:theory} to proposed designs, to check that a design the invertible operator suggests still meets its target once small, realistic manufacturing tolerances are taken into account. Fourth, validating a proposed design experimentally on a real plate, to check that the gap between the solver and reality does not undo what the neural operators learned from simulation alone.

Taken together, this thesis provides forward surrogates for two and three resonators, invertible architectures that already perform generative design with reasonable spectrum and tuning-frequency accuracy, and a reproducible, single-script pipeline for all of it, while identifying position recovery and retraining on the current design space as the concrete next steps toward a complete generative design tool.
